\begin{exercise}{70}
Show that the slice/join adjunction induces a bijection of lifting problems
between diagrams of the kind
\[
 \begin{tikzcd}[column sep=large,row sep=large]
  S \arrow[r] \arrow[d] & X_{\varphi/} \arrow[d] \\
  T \arrow[r] \arrow[ur,dashed]
    & X_{\varphi i/}\times_{Y_{f\varphi i/}}Y_{f\varphi/}
 \end{tikzcd}
\]
and diagrams of the kind
\[
 \begin{tikzcd}[column sep=large,row sep=large]
  B \arrow[r] \arrow[rr,bend left=25,"\varphi"]
    & A\star T\amalg_{A\star S}B\star S
      \arrow[r] \arrow[d]
    & X \arrow[d] \\
    & B\star T \arrow[r] \arrow[ur,dashed]
    & Y
 \end{tikzcd}
\]

\begin{proof}
Let $i\colon A\to B$, $j\colon S\to T$,
$\varphi\colon B\to X$, and $f\colon X\to Y$ be the
maps involved; $i$ and $j$ are independent.
Write $P=A\star T\amalg_{A\star S}B\star S$ and let
$q\colon P\to B\star T$ be induced by $i\star T$ and
$B\star j$.
For $p\colon R\to X$, Corollary~1.4.17 gives
\[
 \mathrm{Hom}(N,X_{p/})
 \cong\{h\colon R\star N\to X\mid h|_R=p\}.
\]
Under this bijection, restriction in $N$, restriction
along $i$, and postcomposition with $f$ become
precomposition with the corresponding join maps and
postcomposition with $f$, respectively
(Lemma~1.4.20).

Transpose the top map of the first square and the two
components of its bottom map. They give
\[
 a\colon B\star S\to X,\qquad
 b\colon A\star T\to X,\qquad
 c\colon B\star T\to Y,
\]
with $a|_B=\varphi$, $b|_A=\varphi i$, and
$c|_B=f\varphi$.
The square and the fiber-product condition say exactly
\[
\begin{aligned}
 a(i\star S)&=b(A\star j),\\
 fa&=c(B\star j),\\
 fb&=c(i\star T).
\end{aligned}
\]
The first equality glues $b$ and $a$ to a unique map
$g\colon P\to X$; the other two give $fg=cq$.
The restriction $a|_B=\varphi$ is precisely the curved
arrow condition in the second diagram.
Conversely, restricting such a diagram to the two
pushout summands recovers $a,b,c$ and these equalities.
The curved arrow fixes $a|_B=\varphi$; gluing then forces
$b|_A=\varphi i$, and commutativity forces
$c|_B=f\varphi$.
Thus transposing back recovers the first square, and
the two constructions are inverse.

Finally, a map $T\to X_{\varphi/}$ transposes to
$h\colon B\star T\to X$ under $B$.
Its three lifting conditions transpose to
\[
 h(B\star j)=a,\qquad h(i\star T)=b,\qquad fh=c.
\]
These are equivalent to $hq=g$ and $fh=c$, exactly the
second dashed lift conditions.
Any such $h$ automatically restricts to $\varphi$ on
$B$, because $B\to B\star T$ factors through
$B\to B\star S\to P\xrightarrow{q}B\star T$.
Hence no extra condition is imposed, and the adjunction
also gives mutually inverse bijections of the dashed
lift sets.
\end{proof}
\end{exercise}
